\section{Introduction}
\label{sec:introduction}

For a complex manifold $M$, let $H^\infty(M)$ denote the algebra of
bounded holomorphic functions on $M$. We say that $M$ has the
\emph{holomorphic Liouville property} when $H^\infty(M)=\mathbb C$.
The question studied here asks whether negative curvature forces this
algebra to contain a nonconstant function. More precisely, the one-sided
bounded-holomorphic-function question asks whether every complete simply
connected K\"ahler manifold of complex dimension at least two satisfying
\begin{equation}\label{eq:one-sided-question}
 \operatorname{Sec}_g\le -a^2<0
\end{equation}
admits a nonconstant bounded holomorphic function. Sectional curvature
here and throughout the paper means curvature on every real two-plane.
There is no lower curvature bound in this hypothesis. We use the public
formulation recorded by Wu and Yau \citep[p.~103]{WuYau2020}; their
survey \citep[Section~3, p.~9]{WuYauSurvey2019} traces it to Yau's
1982 Problem~38.

\begin{theorem}\label{thm:main}
For some finite integer $m\ge2$ there is a domain
$\mathcal T\subset\mathbb C^m$, diffeomorphic to $\mathbb R^{2m}$,
with a complete K\"ahler metric $g$ such that
\[
 \operatorname{Sec}_g\le-1,
 \qquad H^\infty(\mathcal T)=\mathbb C.
\]
The sectional curvatures of $g$ are unbounded below.
\end{theorem}

The dimension is chosen sufficiently large and then fixed for the
entire construction. The theorem therefore disproves the universal
assertion in the one-sided question. Its curvature is unbounded below,
so the example leaves separate the bounded-function question under a
finite two-sided bound
$-b^2\le\operatorname{Sec}_g\le-a^2<0$.

The uniform upper bound is essential to the question. Seshadri,
adapting Klembeck's radial curvature calculations, exhibited a complete
K\"ahler metric on $\mathbb C^n$ with strictly negative sectional
curvature \citep[Section~3]{Seshadri2006}. Since every bounded entire
function is constant, strict negativity by itself already permits the
holomorphic Liouville property. In that example some sectional
curvatures tend to zero; Theorem~\ref{thm:main} retains a uniform
negative upper bound.

\subsection{Bounded functions, bounded coordinates, and invariant metrics}

A biholomorphism to a bounded domain in $\mathbb C^m$ supplies $m$
bounded holomorphic coordinate functions whose differentials are
independent everywhere. More generally, a bounded holomorphic map
$M\to\mathbb C^m$ with nowhere-vanishing Jacobian provides such a tuple,
without necessarily being injective. Either conclusion implies the
existence of a nonconstant bounded holomorphic function. The reverse
implications require additional information.

The stronger bounded-domain question appears in H.~Wu's work
\citep[p.~195, question~(1)]{WuNormal1967}. Its hypotheses are
completeness, simple connectivity, nonpositive real sectional curvature,
and a uniformly negative upper bound for holomorphic sectional
curvature. Theorem~\ref{thm:main} also gives a negative answer under
those hypotheses, since a negative upper bound on all real two-planes
applies in particular to complex lines. Wu and Yau formulate a
bounded-pseudoconvex-domain question under two-sided negative real
sectional pinching in \citep[Conjecture~4.3]{WuYauSurvey2019}.

Comparisons among intrinsic metrics give another approach to these
questions. For complete simply connected K\"ahler manifolds with
two-sided negative real sectional pinching, Greene and Wu established
a complete Bergman metric and a lower comparison with the original
metric. Wu and Yau recall that result and prove the complementary upper
bound \citep[Theorems~4 and~6]{WuYau2020}. Under the same hypotheses,
\citet[Corollary~7]{WuYau2020} compare the Bergman,
Kobayashi--Royden, and complete negative K\"ahler--Einstein metrics with
the original metric. The Bergman metric uses square-integrable
holomorphic top forms, and the Kobayashi--Royden metric uses holomorphic
disks mapped into the manifold. After constant rescaling, bounded
holomorphic functions give maps from the manifold to the disk.

These latter maps define the Carath\'eodory infinitesimal pseudometric:
at each tangent vector, take the supremum of its lengths pulled back
from the Poincar\'e disk metric. Every such map on $\mathcal T$ is constant by
Theorem~\ref{thm:main}, so this pseudometric vanishes identically.
The invariant-metric comparisons above contain no Carath\'eodory
nondegeneracy assertion; they also require the lower curvature bound
absent from our example.

Shcherbina and Zhang construct a domain in $\mathbb C^2$ with smooth
strictly pseudoconvex boundary that is Kobayashi and Bergman complete
and has the holomorphic Liouville property
\citep[Main Theorem and Section~6]{ShcherbinaZhang2021}. Their argument
first forces bounded continuous plurisubharmonic functions to be constant on a
Wermer-type set and then uses holomorphic uniqueness along slices.
Here polynomial sampling and phases enforce bounded-function rigidity
while the construction also controls sectional curvature on every
real two-plane. The announced pinched result of Cao and Shaw
was withdrawn \citep{CaoShawWithdrawn}; it is not an input to our proof.

The separate companion \emph{A negatively pinched K\"ahler threefold
without bounded holomorphic coordinates}
\citep[Theorem~1.1]{OpenAIPinchedThreefold2026} constructs a contractible
domain in $\mathbb C^3$ with a complete two-sided negatively pinched
K\"ahler metric and no bounded holomorphic map to $\mathbb C^3$ with
everywhere nonzero Jacobian. That domain retains two nonconstant bounded
base-coordinate functions. Its conclusion concerns bounded coordinate
tuples, whereas Theorem~\ref{thm:main} concerns every individual bounded
holomorphic function. The two proofs are independent; no result of the
companion is used below.

\subsection{The construction and its main obstacles}

We first separate curvature from bounded-function rigidity. Starting with
a smooth strictly plurisubharmonic potential $\ell$ on
$\mathbb C^{n}_s\times\mathbb C_y$, consider
\begin{equation}\label{eq:intro-tube}
 \mathcal T=\{(s,y,w):\operatorname{Im}w>\ell(s,y)\},
 \qquad g=\partial\bar\partial
       [-\log(\operatorname{Im}w-\ell(s,y))].
\end{equation}
The final complex dimension is $m=n+2$: there are $n$ spatial variables,
one fiber variable $y$, and the tube variable $w$.
If the potential metric of $\ell$ is complete and has nonpositive real
sectional curvature, then the tube metric is complete and satisfies
$\operatorname{Sec}_g\le-1$. Lemma~\ref{tube:transfer} proves this by
comparing the curvature with that of complex hyperbolic space. Thus
we must construct a complete nonpositively curved potential metric
while placing selected large holomorphic disks in its tube.

The disks detect derivatives of bounded holomorphic functions. They
lie over finite sets of spatial points obtained from a homogeneous
polynomial system $(H_N,p_N)$ of degree $N$, where $H_N$ has $n-1$
components and $p_N$ is scalar. The grid consists of selected common
zeros of $H_N$ in projective space $\mathbb P^{n-1}$ at which $p_N$
is nonzero. We use the standard unitary-invariant
Gaussian polynomial ensemble and its determinant-weighted zero-density
formulas, developed by Bleher, Shiffman, and Zelditch
\citep[Theorem~2.1 and Sections~2.2 and~2.4]{BleherShiffmanZelditch2000}.
The distribution of random zeros was studied by Shiffman and Zelditch
\citep{ShiffmanZelditch1999}; Shiffman's later simultaneous-zero theorem
\citep[Corollary~1.3]{Shiffman2008Convergence} includes equidistribution
of the unpruned grids used here. We prove the needed special cases
locally and add quantitative derivative, separation, and lift estimates.

There are on the order of $N^{n-1}$ projective grid points. An affine
lift of such a point is a representative $v\in\mathbb C^n$ normalized
by $p_N(v)=1$. Its $N$ lifts are $e^{2\pi iq/N}v$, $0\le q<N$, and
form a root-of-unity cycle. Averaging around a cycle separates Taylor degrees modulo $N$. In a sufficiently large
fixed dimension, the evaluation vectors of the first few surviving
homogeneous terms span only a small fraction of the space of values
at the grid points. Independent phases assigned to the cycles then
keep a prescribed low-degree jet uniformly away from that subspace.
Each testing disk keeps $s$ fixed while varying $y$ and $w$ with a
phase-dependent slope. Cauchy's estimate on the large disks turns this separation into
exponential decay of the Taylor coefficients of $f_w=\partial f/\partial w$
for every bounded holomorphic $f$. Repeating each jet test at heights that grow slowly
relative to the decay exponent, and applying the three-lines theorem,
forces $f_w=0$. The function then descends to a bounded entire function
on $\mathbb C^{n+1}$ and is constant. Proposition~\ref{grid:criterion}
states the exact disk conditions separately from their realization.

The geometric obstacle is to produce those large disks without losing
curvature control when their centers and directions vary. For the
radial fiber we use a partial Legendre transform, in the circle-invariant
bundle framework of Calabi \citep[Section~3]{Calabi1979} and the momentum
formulation of Hwang and Singer \citep[Section~2.1]{HwangSinger2002}.
Our profile varies over the spatial base; the additional spatial-derivative
terms must be controlled for every real two-plane. The Legendre transform
exposes six curvature blocks. We test their contraction on Hermitian
matrices of rank at most two; every real two-plane produces a matrix
in this class. This rank restriction keeps all trace losses independent of dimension.
A steep fiber profile supplies positive margins in the negative
curvature tensor. They absorb the spatial polynomial wells and small
nonholomorphic shear coefficients, from the fiber axis to infinity.

Successive tests require different fiber centers. A direct center
motion is not a holomorphic coordinate change and can destroy the
near-axis estimates. We align the jets of two radial spatial metrics,
temporarily widen the fiber profile, move the center while this width
is available, and restore the profile with a controlled constant
offset. Each period of rapid spatial growth occupies a finite interval
and is followed by a capped growth regime, so no finite spatial radius
is lost to blowup. Finally, every stage reserves bounds for all future
wells. The local estimates therefore apply to the final potential on
entire fibers and preserve every previously installed disk.

Sections~\ref{sec:tube} and \ref{sec:analytic} establish the curvature
transfer and the analytic testing criterion. Section~\ref{sec:fiber}
derives the Legendre metric and curvature formulas.
Section~\ref{sec:quiet} gives the profile and shear estimates;
Sections~\ref{sec:radial} and~\ref{sec:centers} supply the aligned spatial
geometry and center motion. Section~\ref{sec:assembly} fixes the
constants before choosing the dimension, carries out the global
recursion, proves completeness, and concludes Theorem~\ref{thm:main}.

\subsection{Conventions and uniformity}

We measure squared lengths using the Hermitian metric matrix:
$|v|_g^2=g_{i\bar j}v^i\overline{v^j}$ for the complex components of a
real tangent vector. In particular a potential $\ell$ has metric matrix
$(\ell_{i\bar j})$. This convention gives sectional curvatures between
$-4$ and $-1$ for the ball potential $-\log(1-|z|^2)$.
For the same complex Hessian, the companion
\citep[Section~2]{OpenAIPinchedThreefold2026} uses twice this real metric;
its complex-hyperbolic sectional-curvature interval is therefore
$[-2,-1/2]$.
The negative Hermitian curvature tensor is denoted by $C=-R$; its
coordinate formula and relation to real sectional curvature are given
in Section~\ref{sec:fiber}.

The symbols $O(\cdot)$, $\lesssim$, and $\asymp$ denote respectively
a constant-factor bound, an upper bound, and two-sided positive bounds.
Dependencies are stated when the bound is introduced. Tensor jets are
measured in fixed affine coordinates at the point, using operator or
multilinear norms, rather than sums of their entries. Hermitian test
matrices carry the Frobenius norm. For a rank at most two test $T$ and a
Hermitian form $Q$ this gives
\[
 |Q:T|\le\sqrt2\,\|Q\|\,|T|.
\]
The constants that determine the analytic exponent gaps and the eventual
dimension are independent of dimension. Constants used after the
dimension has been fixed may depend on it. The order of choices is
specified in Section~\ref{sec:assembly}; every choice respects that
order of dependence.
